% Generated draft: requires mathematical and source review.
\documentclass[11pt]{article}
\usepackage[margin=25mm]{geometry}
\usepackage{fontspec}
\setmainfont{Latin Modern Roman}
\usepackage{amsmath,amssymb,mathtools,unicode-math}
\setmathfont{Latin Modern Math}
\usepackage{enumitem,needspace,xurl,hyperref}
\setlength{\parindent}{0pt}
\setlength{\parskip}{5pt}
\setlength{\emergencystretch}{3em}
\newcommand{\N}{\mathbb{N}}
\newcommand{\Z}{\mathbb{Z}}
\newcommand{\Q}{\mathbb{Q}}
\newcommand{\R}{\mathbb{R}}
\newcommand{\C}{\mathbb{C}}
\newcommand{\F}{\mathbb{F}}
\newcommand{\A}{\mathbb{A}}
\newcommand{\PP}{\mathbb{P}}
\newcommand{\E}{\mathbb{E}}
\newcommand{\eps}{\varepsilon}
\newcommand{\dd}{\,\mathrm{d}}
\begin{document}
\section*{UnsolvedMath ID: 20002074}
\textbf{Source problem code:} catalogue entry
\subsection{Definitions and mathematical statement}
\textbf{Question 1.16, Auroux and workshop discussion.} Construct a category of Legendrians in a contact open book analogous to a Fukaya category, retaining mixed Reeb chords, self-contact-homology coefficient data, and winding around the open-book angle. For Legendrians associated to one page, the source warns that a single Lagrangian Floer complex omits iterated monodromy contributions. It further asks whether restoring the binding can be described as an $A_\infty$ deformation, and suggests that relative contact homology should be a module over absolute contact homology. The exact contact-homology theory, objects and analytic hypotheses are unspecified.

\textbf{Supplied algebraic orbit model.} Let $\mathcal A$ be a strictly unital uncurved $A_\infty$ category over a field $\mathbb k$, and let $\Phi$ be a strict $A_\infty$ automorphism preserving units and all operations. Define a nonnegative orbit category with the same objects and
\[
 \operatorname{hom}_{\mathcal O_+}(X,Y)
    =\bigoplus_{m\ge0}\operatorname{hom}_{\mathcal A}(X,\Phi^mY)t^m,
\]
where $t$ records winding and is assigned cohomological degree zero in this model. For composable $f_i\in\operatorname{hom}_{\mathcal A}(X_{i-1},\Phi^{m_i}X_i)$, put $M_j=m_1+\cdots+m_j$ and define
\[
 \mu_{\mathcal O_+}^d(f_dt^{m_d},\ldots,f_1t^{m_1})
  =\mu_{\mathcal A}^d(\Phi^{M_{d-1}}f_d,\ldots,
                         \Phi^{M_1}f_2,f_1)t^{M_d}.
\]
For $d=1$ the formula is simply $\mu^1(f_1)t^{m_1}$. The aid asserts the $A_\infty$ identities and strict units at winding zero. When $\Phi$ is the identity, this is $\mathcal A\otimes_{\mathbb k}\mathbb k[t]$, retaining all nonnegative iterates rather than one copy of a Floer complex. Mixed morphisms carry the canonical higher bimodule structure over the endomorphism $A_\infty$ algebras at their two endpoints; replacing these by bare homology groups can lose the required coefficients and operations.

\textbf{Supplied winding-defect deformation.} Let the bar coderivation have a genuine homogeneous nonnegative winding-defect decomposition $b=\sum_{m\ge0}b_m$, with degree-one coderivations $b_m$ and additive defect under composition. Equivalently the homogeneous $A_\infty$ identities give
\[
 \sum_{r+s=m}b_rb_s=0\quad\text{for every }m\ge0.
\]
Assume local finiteness on each input. Then the aid's Rees family
\[
 b_q=\sum_{m\ge0}q^m b_m
\]
satisfies $b_q^2=0$, with $q=0$ the defect-zero structure and $q=1$ the original one. The index on $b_m$ denotes defect, not operation arity or a curvature term. If local finiteness is dropped, a completed formal family may require additional convergence data, and specialization at $q=1$ is not automatic.

The proposed geometric application is conditional: holomorphic-curve counts for the open-book complement and restored binding must furnish these winding pieces, compactness, regularity, boundary-of-moduli-space identities, and gluing identifications. Actual monodromy may act only up to coherent homotopy, while the displayed model assumes a strict automorphism. Neither strictification nor the geometric equivalence is proved in the supplied record. Nonnegative winding is part of this model and is not imposed without qualification on every possible open-book contact-homology construction.

Eliashberg's source discussion describes convex, flat and concave boundary components in the symplectization of the binding complement, rounding corners to absorb the flat part into the convex part, and partial gluing along that convex boundary when the binding is restored. This motivates the deformation proposal but does not give the missing analytic theorem. The trailing ``Third day'' and Paul Biran introduction begin subsequent questions that are absent here.
\subsection{Short English statement}
Auroux asks for a Legendrian category in an open book that remembers self-algebras, mixed chords and repeated monodromy. The supplied algebraic model has a positive orbit category and a winding Rees deformation; identity monodromy gives a polynomial family of morphisms. Identifying this model with binding restoration remains conditional on geometric counts, compactness and gluing.
\subsection{Sources}
\begin{itemize}
\item Local UnsolvedMath registry, permanent ID 20002074, source code catalogue entry; canonical \texttt{aim-geometry-notes.json}, notes[411].
\item AIM \emph{Holomorphic curves in contact geometry}, Question 1.16, Auroux with workshop discussion and Eliashberg's gluing comments: \url{https://aimath.org/WWN/contactgeom2/contactgeom2.pdf}. The supplied verification reports comparison with PDF pp.~5--6, not performed here.
\item No separate bibliographic citation is supplied for the orbit-category construction, Rees-family reduction or the required contact-geometric realization theorem. The Paul Biran line introduces later material, not this question's attribution.
\item Registry research attempt 1, dated 2026-08-14: reduction, rigor ``proved'', novelty confidence low. The machine-generated, non-peer-reviewed aid requires independent verification before citation.
\end{itemize}
\end{document}
